\documentclass[11pt,a4paper]{article}

\usepackage[utf8]{inputenc}
\usepackage[T1]{fontenc}
\usepackage[english]{babel}
\usepackage{xcolor}
\usepackage{amsmath,amssymb,amsthm}
\usepackage{amscd}
\usepackage{mathtools}
\usepackage[expansion=false]{microtype}
\setlength{\emergencystretch}{2em}
\usepackage{geometry}
\geometry{margin=2.4cm}
\usepackage{booktabs}
\usepackage{tabularx}
\usepackage{array}
\usepackage{ragged2e}
\newcolumntype{Y}{>{\RaggedRight\arraybackslash}X}
\usepackage{enumitem}
\usepackage[most]{tcolorbox}
\usepackage{hyperref}
\hypersetup{
  colorlinks=true,
  linkcolor=blue!60!black,
  citecolor=blue!60!black,
  urlcolor=blue!60!black
}

% theorem environments
\newtheorem{theorem}{Theorem}[section]
\newtheorem{proposition}[theorem]{Proposition}
\newtheorem{corollary}[theorem]{Corollary}
\newtheorem{lemma}[theorem]{Lemma}
\newtheorem{conjecture}[theorem]{Conjecture}
\theoremstyle{definition}
\newtheorem{definition}[theorem]{Definition}
\newtheorem{postulate}{Postulate}
\newtheorem{hypothesis}{Hypothesis}
\theoremstyle{remark}
\newtheorem{remark}[theorem]{Remark}
\newtheorem{anchor}[theorem]{Physical anchor}

% shortcuts
\newcommand{\Sn}[1]{S^{#1}}
\newcommand{\Szero}{S^{0}}
\newcommand{\Sone}{S^{1}}
\newcommand{\Stwo}{S^{2}}
\newcommand{\Sthree}{S^{3}}
\newcommand{\Sfour}{S^{4}}
\newcommand{\Sseven}{S^{7}}
\newcommand{\R}{\mathbb{R}}
\newcommand{\C}{\mathbb{C}}
\newcommand{\Z}{\mathbb{Z}}
\newcommand{\Hbb}{\mathbb{H}}
\newcommand{\CP}{\mathbb{C}P}
\newcommand{\SU}{\mathrm{SU}}
\newcommand{\U}{\mathrm{U}}
\newcommand{\lP}{\ell_{P}}
\newcommand{\tP}{t_{P}}
\newcommand{\rH}{\widetilde H}
\newcommand{\Dd}{\mathsf{D}}
\newcommand{\Dsharp}{\mathsf{D}^{\sharp}}
\newcommand{\Dtop}{\mathsf{D}^{\sharp}_{\mathrm{top}}}
\newcommand{\Dop}{\mathsf{D}^{\sharp}_{\mathrm{op}}}
\newcommand{\Dchar}{\mathsf{D}^{\sharp}_{\mathrm{char}}}
\newcommand{\Dspec}{\mathsf{D}^{\sharp}_{\mathrm{spec}}}
\newcommand{\Car}{\mathbf{Car}}
\newcommand{\Str}{\mathbf{Str}}
\newcommand{\Forg}{\mathbf{Forg}}
\newcommand{\Ho}{\mathrm{Ho}}
\newcommand{\golden}{\varphi}
\newcommand{\PIR}{\mathrm{PIR}}
\newcommand{\BCR}{\mathrm{BCR}}
\newcommand{\nkin}{n_{\mathrm{kin}}}
\newcommand{\ncausal}{n_{\mathrm{causal}}}
\newcommand{\Bphys}{B_{\mathrm{phys}}}
\newcommand{\Bcausal}{B_{\mathrm{causal}}}
\newcommand{\Bmode}{b_{\mathrm{mode}}}
\newcommand{\tcoh}{t_{\mathrm{coh}}}
\newcommand{\Rm}{\mathrm{Rm}}

\title{\bfseries Dependently Reconstructible Distinguishability (DRD):\\
A Functorial Theory of Information Carriers,\\
Graded Spectral Capacities, and Causal Integration\\
on Compact \texorpdfstring{$\Sthree$}{S3} Geometry}

\author{%
\texorpdfstring{$^{1}$}{1}Andrei Preece\thanks{\texttt{andre.preece@web.de}; ORCID: 0009-0004-5978-795X}
\quad and \quad
\texorpdfstring{$^{1}$}{1}Boris Batenin\thanks{\texttt{batenin@web.de}; ORCID: 0009-0000-7027-8583}\\[0.35em]
\small \texorpdfstring{$^{1}$}{1}Independent Researchers
}

\date{Berlin, 2026}

\begin{document}
\maketitle

\begin{abstract}
Existing quantitative theories of information are measure theories over presupposed alphabets:
Shannon's theory presupposes a symbol set, algorithmic complexity a machine model, quantum
information a Hilbert space.  This paper develops a definition-first framework---\emph{dependently
reconstructible distinguishability} (DRD)---for the layer beneath those theories: what a
distinction is, what carries it, which operations preserve it, which destroy it and at what
price, and how much of it a physical domain can carry and causally integrate.

The framework is organized in three strictly separated layers.  (I)~A \emph{mathematical layer}:
the correct object is not one flat functor but a fibrewise profile over a category of structures,
$\Car_*^{\mathrm{fib}}=\int_{\mathfrak s\in\Str}\Car_{\mathfrak s}$, with
$\Dsharp_{\mathfrak s}:\Car_{\mathfrak s}\to\mathbf{Prof}_{\mathfrak s}$.  Homology,
cohomology, characteristic classes, spin data and spectra are preserved only inside the
structural fibre whose morphisms preserve the relevant structure.  (II)~An \emph{erasure and
capacity layer}: non-invertible change is measured by a cone-defect package
$\mathbb E(f)$, sector erasure has the safe Landauer component $k_BT\ln N$ plus
platform-dependent geometric cost, and the graded scalar, TT and Dirac capacities on round
$\Sthree_R$ are counted with explicit dimension and cutoff conventions.  Heat-trace capacity
$\mathcal N_P(s,\tau)=\operatorname{Tr}(e^{-sP_\tau})$ replaces discontinuous sharp counts
when a reconstruction flow varies the metric.  (III)~A \emph{physical-selection layer}: causal
integration is exact for shell-complete homogeneous spectral projectors, while the only central
remaining physical node is the gravitational selector.  It is not a ratio between unlike
quantities: the raw spectral count $N_P$ is a mode number, while the selected physical object is
$\log_2\dim\mathcal H_{\mathrm{code}}$ in bits.  The selector is therefore defined only after
an operational finite-alphabet or energy/occupation cutoff has converted accessible modes into
accessible code bits.  It then chooses the area-compatible, gauge-invariant,
backreaction-stable code subspace from the volume-scaling kinematic alphabet.  The framework does not replace Shannon, Holevo, or Landauer; it supplies
the pre-statistical object---the physically realizable alphabet of distinguishable states---that
those theories presuppose.
\medskip

\noindent\textbf{Keywords:} theory of information; distinguishability; carriers; functorial
invariants; suspension; $\Sthree$; Hopf holonomy; spectral capacity; causal integration;
Landauer principle; covariant entropy bounds; boundary admissibility.
\end{abstract}

\tableofcontents

\section{Introduction}

\subsection{The problem: information named before it is defined}

The word ``information'' now carries load in thermodynamics, quantum theory, cosmology, and the
foundations of computation.  Yet the standard quantitative theories are measure theories over
\emph{presupposed} alphabets.  Shannon's theory~\cite{Shannon1948} quantifies the statistics of
a source once a set of distinguishable symbols is given; Kolmogorov
complexity~\cite{Kolmogorov1965} quantifies description length once a machine model and a
symbol set are given; quantum information theory quantifies state distinguishability once a
Hilbert space is given, with the Holevo bound~\cite{Holevo1973} governing accessible classical
information.  None of them answers the prior questions: what is a distinction, what does it
take for a distinction to be physically realized, what operations preserve it, what operations
destroy it and at what price, and how much distinguishability a given physical domain can carry
at a given resolution---and causally integrate in a given time.

Foundational slogans have tried to close the gap from the other side.  Wheeler's ``it from
bit''~\cite{Wheeler1990} makes the bit primitive and matter derivative; Landauer's
``information is physical''~\cite{Landauer1991} makes the carrier primitive.  Both intuitions
are correct in part, and both are incomplete without a formal object that says what a ``bit''
is before any physics, and what physics must add before that object becomes operational.

This paper supplies such an object and develops it as a research framework with an explicit
claim-status discipline.  In its strict form, DRD is best understood not as ``a theory of
information in general'' but as a \textbf{theory of invariants of reconstructible
distinguishability on categories of carriers}.  Information is not identified with statistics,
with entropy, or with a quantum state; it is defined as \emph{structure-preserved
distinguishability}, which acquires physical realization only after metric, causal order, spin
structure, a resolution scale, and a class of admissible morphisms are added.  This contradicts
neither Shannon, nor Holevo, nor Landauer; it places them on different floors of one tower.

The paper proceeds in three layers:

\begin{enumerate}[label=(\Roman*)]
\item a \textbf{mathematical layer} (Secs.~\ref{sec:profile}--\ref{sec:theorems}): the
fibred carrier category $\Car_*^{\mathrm{fib}}$, the fibrewise profiles
$\Dsharp_{\mathfrak s}$, and the conservation, attributed-erasure, and transport theorems;
\item an \textbf{operational layer} (Secs.~\ref{sec:physical}--\ref{sec:pir}): two explicit
postulates, exact graded capacities for scalar, tensor, and spinor alphabets on
$\Sthree_{R}$, and the causal-integration formula $\PIR(T)$;
\item an \textbf{interpretive and programmatic layer}
(Secs.~\ref{sec:qk}--\ref{sec:interpretation}), explicitly labelled: erasure-cost hypotheses
$Q_{k}$, carrier-resolved quantum resources, the volume--area problem, related definition-first
programmes, and the ontological reading.
\end{enumerate}

The companion papers of this programme supply geometric infrastructure: the suspension-chain
analysis of $\Sthree$ as minimal spin/Hopf carrier~\cite{PB2026susp}, the Ricci-type
reconstruction flow with its no-go theorems and Planck-barrier invariant
region~\cite{BP2026ricci}, and the compact causal accessibility structure of light on
$\Sthree_{R}$~\cite{PB2026light}.  The present paper is the information-theoretic keystone of
that line: it states what the other papers' carriers carry.

\subsection{Claim-status convention}

\begin{center}
\small
\begin{tabularx}{\textwidth}{@{}lY@{}}
\toprule
\textbf{Label} & \textbf{Meaning} \\
\midrule
Theorem / Proposition / Lemma & Proved under stated mathematical assumptions.\\
Physical anchor & Established physics cited for contact, not re-derived here.\\
Postulate & An operational assumption of the physical layer, explicitly adjustable.\\
Hypothesis & A proposed law-like statement with specified test platforms; not a theorem.\\
Ansatz & A proposed quantitative form, not derived from dynamics.\\
Interpretation & Philosophy-of-physics reading; proves no physics.\\
\bottomrule
\end{tabularx}
\end{center}

\subsection{What this paper does not claim}

\begin{tcolorbox}[colback=gray!4,colframe=gray!50,title=Boundary of the paper]
This paper does not replace or amend Shannon's theory, Kolmogorov complexity, the Holevo bound,
or Landauer's principle; it supplies the pre-statistical object those theories presuppose.  It
does not derive physical law from topology; does not derive the value of $\lP$; does not prove
that the observed universe is globally $\Sthree$; does not derive gravitational entropy bounds;
does not resolve the volume--area tension (it states it as the central open problem); and does
not assert any ontological statement as a theorem.  Every formula in the operational layer is
conditional on the stated postulates.  Statements of the form ``information is neither created
nor destroyed'' are made only in the precise form: \emph{within a fixed category of invertible
structure-preserving morphisms, the corresponding invariant package is neither created nor
erased}.  The remaining physical problem is explicitly localized in the selector
$\Sigma_{\mathrm{grav}}$, not hidden inside topological language.
\end{tcolorbox}

\begin{theorem}[The DRD closure map]
\label{thm:main-pipeline}
Under the structural and operational assumptions stated in this paper, the usable information
capacity of a sector is not a primitive number but the output of four separated maps:
\[
\Dsharp_{q,\Lambda}(\mathcal X)
\;\rightsquigarrow\;
\nkin^{(P)}
\;\xrightarrow{\ \PIR_P(T)\ }\;
\ncausal^{(P)}
\;\xrightarrow{\ \Bmode^{(P)}\ }\;
\Bcausal^{(P)}
\;\xrightarrow{\ \Sigma_{\mathrm{grav}}^{(P)}\ }\;
\Bphys^{(P)}.
\]
The first object is the fibrewise distinguishability profile; $\nkin^{(P)}$ and
$\ncausal^{(P)}$ are mode numbers; $\Bcausal^{(P)}$ and $\Bphys^{(P)}$ are bit counts after a
finite alphabet, energy cutoff or code-space cutoff has been specified.  The mathematical and
spectral parts of the map are computable inside their declared fibres.  The only irreducible
physical selector left open is $\Sigma_{\mathrm{grav}}^{(P)}$, the map from accessible code bits
to gauge-invariant, gravitationally dressed, entropy-admissible, backreaction-stable code bits.
\end{theorem}

\begin{remark}[What the closure theorem forbids]
Equation~\eqref{thm:main-pipeline} is a bookkeeping discipline.  It forbids three common
mistakes: treating a vector-space dimension as a bit count, treating a boundary entropy bound as
a replacement of the bulk carrier, and treating a suspension isomorphism in reduced homology as
an isomorphism of the entire unstable profile.
\end{remark}


\section{The Full Distinguishability Profile \texorpdfstring{$\Dsharp$}{D-sharp}}
\label{sec:profile}

\subsection{Elementary distinction}

\begin{definition}[Elementary distinction]
An \emph{elementary distinction} is the two-point space $\Szero=\{-1,+1\}$ with its discrete
topology.  No metric, phase, direction, or causal structure is assumed at this level.  The two
points are not two objects located somewhere; they are the two poles of a single relational
structure, \emph{this / not-this}.  Neither point carries the structure alone: removing either
yields a one-point space, which is no longer a distinction.
\end{definition}

\subsection{Why a single rank profile is not enough}

A first-pass invariant is the rank profile
$\Dd_{k}(X)=\operatorname{rank}\rH_{k}(X;\Z)$: $\Dd_{0}$ counts independent elementary
distinctions, $\Dd_{1}$ cyclic (phase-type) distinctions, $\Dd_{2}$ enclosure-type
distinctions, and so on; for spheres $\Dd_{k}(\Sn n)=\delta_{k,n}$.  The rank profile already
supports invariance and transport theorems, but it conflates levels of structure: it cannot see
torsion, ring structure, unstable homotopy data, bundle topology, or resolution-dependent
spectral content.  A publishable theory must therefore carry a \emph{two-level functorial
package}, not one number.

\begin{definition}[Topological profile]
\label{def:dtop}
For a pointed space $X$ of the homotopy type of a finite CW complex,
\begin{equation}
\Dtop(X)\;:=\;
\Big(
\{\rH_{k}(X;\Z)\}_{k\ge0},\;
H^{\ast}(X;\Z)\ \text{with}\ \smile,\;
\{\pi_{k}(X)\}_{k\ge1}\ \text{with Whitehead products}
\Big),
\end{equation}
where $\rH_{k}(X;\Z)\cong\Z^{\beta_{k}}\oplus\mathrm{Tor}_{k}$ retains torsion.  The rank
shadow is $\Dd_{k}(X)=\beta_{k}$.
\end{definition}

\begin{definition}[Structured carrier and operational profile]
\label{def:dop}
A \emph{structured carrier} is a tuple
\begin{equation}
\mathcal X=(X,g,\mathfrak s,\mathcal E,\varepsilon),
\end{equation}
where $g$ is a Riemannian metric, $\mathfrak s$ a spin structure (when required),
$\mathcal E$ a finite set of distinguished vector bundles/fields, and $\varepsilon>0$ a
resolution scale.  Its \emph{operational profile} is
\begin{equation}
\Dop(\mathcal X)\;:=\;
\Big(
\Dtop(X),\;
\mathcal C(\mathcal X),\;
\mathrm{Spec}(\mathcal X)
\Big),
\end{equation}
where $\mathcal C(\mathcal X)$ is the package of characteristic classes
($w_{i},c_{i},p_{i},e,\hat A$) of the bundles in $\mathcal E$, and
$\mathrm{Spec}(\mathcal X)=(\mathrm{Spec}_{0},\mathrm{Spec}_{TT},\mathrm{Spec}(D_{\mathfrak s}))$
collects the spectra (with multiplicities) of the scalar Laplacian, the rough (Bochner)
Laplacian on transverse-traceless tensors, and the Dirac operator, truncated at the resolution $\varepsilon$ where
an operational alphabet is required.
\end{definition}

\subsection{Functoriality}

\begin{theorem}[Functoriality of the topological profile]
\label{thm:functorial}
A continuous map $f:X\to Y$ induces
$f_{\ast}:\rH_{k}(X;\Z)\to\rH_{k}(Y;\Z)$ (covariant),
$f^{\ast}:H^{\ast}(Y;\Z)\to H^{\ast}(X;\Z)$ (contravariant) with
$f^{\ast}(\alpha\smile\beta)=f^{\ast}\alpha\smile f^{\ast}\beta$, and
$\pi_{k}(f):\pi_{k}(X)\to\pi_{k}(Y)$ compatible with Whitehead products.  Characteristic
classes are natural under pullback: $x(f^{\ast}E)=f^{\ast}x(E)$.  Hence $\Dtop$, and the class
package $\mathcal C$, are functorial with respect to (admissible) continuous maps.
\end{theorem}

\begin{proof}
Standard algebraic topology and the theory of characteristic
classes~\cite{Hatcher2002,MilnorStasheff}.
\end{proof}

\begin{proposition}[Spectral invariance under (spin-)isometry]
\label{prop:spec}
If $f:(X,g)\to(Y,h)$ is an isometry, the scalar and TT-tensor Laplace-type operators are unitarily
conjugate, so $\mathrm{Spec}_{0}$ and $\mathrm{Spec}_{TT}$ coincide with multiplicities; if $f$
is in addition a spin-isometry, $\mathrm{Spec}(D_{\mathfrak s})$ coincides as well.  Thus
$\mathrm{Spec}$ is functorial on the subcategory of (spin-)isometric morphisms, not on all
continuous maps.
\end{proposition}

\begin{proof}
Pullback by an isometry intertwines the Laplace-type operators; a spin-isometry lifts to a
unitary on spinor fields intertwining the Dirac operators~\cite{Baer1996}.
\end{proof}

\begin{remark}[The point of the two-level structure]
With Definitions~\ref{def:dtop}--\ref{def:dop}, DRD is \emph{not one number but a functorial
package of distinguishability}, whose components live at different structural levels and
degrade differently under different forgetting operations.  This removes the ambiguity of
identifying ``information'' sometimes with a homology rank and sometimes with a spectral
capacity: after the split into $\Dtop$ and $\Dop$ these are not in conflict but form a
hierarchy of levels.
\end{remark}

\begin{center}
\small
\begin{tabularx}{\textwidth}{@{}lYYY@{}}
\toprule
\textbf{Layer} & \textbf{Object} & \textbf{Functoriality} & \textbf{What it registers}\\
\midrule
Homological & $\rH_{k}(X;\Z)\cong\Z^{\beta_k}\!\oplus\mathrm{Tor}_k$ & covariant ($f_{\ast}$) &
components, cycles, enclosures, torsion distinctions\\
Cohomological & $H^{\ast}(X;\Z)$, $\smile$ & contravariant ($f^{\ast}$) &
intersection/ring structure, obstruction data\\
Homotopical & $\pi_{k}(X)$, Whitehead products & covariant ($\pi_{k}(f)$) &
unstable distinctions invisible to homology\\
Bundle & $w_{i},c_{i},p_{i},e,\hat A$ of $\mathcal E$ & contravariant (pullback) &
orientation, spin, complex/real field topology\\
Spectral & $\mathrm{Spec}_{0},\mathrm{Spec}_{TT},\mathrm{Spec}(D)$ &
invariant under (spin-)isometries & operationally resolvable capacity\\
\bottomrule
\end{tabularx}
\end{center}

\subsection{Strengthened formulation: the structure fibration}
\label{subsec:structure-fibration}

The most robust way to close the object $\Dsharp$ is not to force all invariants into one
flat category.  Homology, characteristic classes, spin data and spectra have different
admissible morphisms.  The correct ambient object is therefore a small category of structural
packages together with its Grothendieck construction.

\begin{definition}[Structure fibration]
Let $\mathbf{Str}$ be the category whose objects are packages
\[
\mathfrak s=(\mathrm{Top},\mathrm{CW},\mathrm{Smooth},\mathrm{Orient},
\mathrm{Spin/Pin},\mathrm{Bundle},\mathrm{Metric},\mathrm{Operator},\varepsilon,
\Lambda,q),
\]
and whose morphisms forget or weaken structure.  For each $\mathfrak s$ let
$\Car_{\mathfrak s}$ be the category of carriers preserving exactly the structure listed in
$\mathfrak s$.  The full carrier category is the Grothendieck construction
\[
\Car_{*}^{\mathrm{fib}}:=\int_{\mathfrak s\in\mathbf{Str}}\Car_{\mathfrak s}.
\]
A profile is then a family of functors
\[
\Dsharp_{\mathfrak s}:\Car_{\mathfrak s}\longrightarrow \mathbf{Prof}_{\mathfrak s},
\]
compatible with weakening of structure whenever a forgetting morphism is declared.  Concretely,
for a weakening $\alpha:\mathfrak s\to\mathfrak t$ the square
\[
\begin{CD}
\Car_{\mathfrak s} @>{\Dsharp_{\mathfrak s}}>> \mathbf{Prof}_{\mathfrak s}\\
@V{U_\alpha}VV @VV{V_\alpha}V\\
\Car_{\mathfrak t} @>{\Dsharp_{\mathfrak t}}>> \mathbf{Prof}_{\mathfrak t}
\end{CD}
\]
commutes up to the declared comparison isomorphism.  Here $U_\alpha$ forgets carrier structure
and $V_\alpha$ forgets the corresponding profile components.
\end{definition}

\begin{proposition}[Layer-compatibility]
Inside a fixed fibre $\Car_{\mathfrak s}$, isomorphisms preserve every component of
$\Dsharp_{\mathfrak s}$.  Across fibres, a change of profile is not a paradox but an exhibited
weakening of structure.  Thus the full invariant is best written operationally as
\[
\Dsharp_{q,\Lambda}(\mathcal X)=
\big(\mathsf{D}^{\sharp}_{\mathrm{top},\le q}(X),\ \mathcal C(\mathcal X),\
\mathrm{Spec}_{\le\Lambda}(\mathcal X)\big),
\]
with $q$ a homotopy/Postnikov cutoff and $\Lambda\sim\varepsilon^{-1}$ a spectral cutoff.
\end{proposition}

\begin{remark}[Status of the closure]
This closes OP1 as a mathematical construction.  The formal object is the fibrewise profile
$\Dsharp_{\mathfrak s}$; the publishable and computable object is the truncated profile
$\Dsharp_{q,\Lambda}$; and the physical object is the same profile after causal and
gravitational admissibility have been applied.
\end{remark}

\section{Carrier Categories, Forgetting Morphisms, and the Suspension Functor}
\label{sec:categories}

\subsection{Two categories of carriers}

\begin{definition}[Carrier categories]
\label{def:categories}
$\Car^{\mathrm{top}}_{\ast}$ is the category of pointed carriers of finite CW type with
basepoint-preserving continuous maps.  $\Car^{\mathrm{op}}$ is the category of structured
carriers $\mathcal X=(X,g,\mathfrak s,\mathcal E,\varepsilon)$
(Definition~\ref{def:dop}); its morphisms preserve exactly those structures needed to define
the corresponding operators and spectra (isometries for $\mathrm{Spec}_{0},\mathrm{Spec}_{TT}$;
spin-isometries for $\mathrm{Spec}(D)$; bundle maps covering the base map for $\mathcal C$).
\end{definition}

\begin{definition}[Forgetting morphisms]
\label{def:forg}
$\Forg\subset\mathrm{Mor}(\Car)$ is the class generated by four operation types:
\begin{enumerate}[label=(\roman*)]
\item \emph{quotient / collapse} (identification of points or subspaces);
\item \emph{deletion / restriction} (removal of points, subspaces, or components);
\item \emph{change of coefficients / change of category} (e.g.\ passing to a homotopy-coarser
category, discarding torsion, forgetting a bundle or the spin structure);
\item \emph{metric--spectral coarse graining} (increasing the resolution scale
$\varepsilon\mapsto\varepsilon'>\varepsilon$).
\end{enumerate}
\end{definition}

\begin{lemma}[Coarse graining is monotone on capacity]
\label{lem:coarse}
For fixed $(X,g)$ and $\varepsilon'\ge\varepsilon$, the band-limited spaces nest,
$\mathcal A_{\varepsilon'}\subseteq\mathcal A_{\varepsilon}$, so every spectral capacity
satisfies $N_{\varepsilon'}\le N_{\varepsilon}$.
\end{lemma}

\begin{proof}
The eigenvalue condition $\sqrt\lambda\le1/\varepsilon'$ implies $\sqrt\lambda\le1/\varepsilon$.
\end{proof}

\subsection{Suspension as an endofunctor}

\begin{theorem}[Transport; degree-shift conservation]
\label{thm:transport}
The reduced suspension $\Sigma X=\Sone\wedge X$ is an endofunctor
$\Sigma:\Car^{\mathrm{top}}_{\ast}\to\Car^{\mathrm{top}}_{\ast}$, and there is a natural
isomorphism of functors
\begin{equation}
\sigma_{k}:\ \rH_{k}\;\Longrightarrow\;\rH_{k+1}\circ\Sigma,
\qquad
\rH_{k+1}(\Sigma X)\cong\rH_{k}(X)\quad\text{for all }k\ge0 .
\end{equation}
Hence the reduced homological profile is shifted by one degree without loss:
$\Dd_{k+1}(\Sigma X)=\Dd_{k}(X)$, torsion included.  Suspension changes the carrier and the
degree of every homological distinction; it need not preserve the full unstable homotopy profile
outside the stability range specified below.
\end{theorem}

\begin{proof}
Mayer--Vietoris on the two cones covering $\Sigma X$, both contractible, with intersection
homotopy equivalent to $X$; naturality is standard~\cite{Hatcher2002}.
\end{proof}

\begin{theorem}[Freudenthal stabilization of homotopy data]
\label{thm:freudenthal}
If $X$ is an $(n{-}1)$-connected CW complex, the suspension homomorphism
$\pi_{i}(X)\to\pi_{i+1}(\Sigma X)$ is an isomorphism for $i<2n-1$ and surjective for
$i=2n-1$~\cite{Hatcher2002}.  Thus suspension transports not only homological content but also
homotopy-level distinguishability in the stable range, stabilizing the unstable part of
$\Dtop$.
\end{theorem}

\begin{remark}[Stability discipline]
Suspension is used here as a transport of reduced homology and, by
Theorem~\ref{thm:freudenthal}, as a stable-range transport of homotopy.  It is not asserted to
be fully faithful on all pointed homotopy types.  Thus no conclusion of the paper depends on
the false statement that $\Sigma X$ preserves every unstable invariant of $X$.
\end{remark}


\begin{corollary}[The fundamental class of $\Sthree$ is a suspended distinction]
\label{cor:s3class}
Along $\Szero\xrightarrow{\Sigma}\Sone\xrightarrow{\Sigma}\Stwo\xrightarrow{\Sigma}\Sthree$,
\begin{equation}
\rH_{0}(\Szero)\simeq\rH_{1}(\Sone)\simeq\rH_{2}(\Stwo)\simeq\rH_{3}(\Sthree)\simeq\Z ,
\end{equation}
carrying the distinction class of $\Szero$ to a generator of $H_{3}(\Sthree;\Z)$
\cite{PB2026susp}.  The elementary distinction, the phase cycle, the directional enclosure, and
the compact volume class are one invariant on four carriers.
\end{corollary}

\begin{remark}[Enrichment: what higher carriers add]
Suspension preserves content but changes what the carrier can \emph{do} with it.  On $\Sone$
the distinction acquires phase ($\U(1)$); on $\Stwo$, direction ($\CP^{1}$, the Bloch sphere);
on $\Sthree$, simultaneously: compactness, simple connectivity, the non-Abelian group
$\SU(2)$, spin structure, parallelizability, and the nontrivial Hopf bundle
$\Sone\hookrightarrow\Sthree\to\Stwo$ with $c_{1}=1$.  Within the axiomatics of the companion
paper, $\Sthree$ is the \emph{minimal joint carrier of this stated package of
structures}~\cite{PB2026susp}---a minimality claim relative to the suspension chain and the
listed conditions, not an absolute fundamentality claim.
\end{remark}

\section{Conservation and Attributed Erasure}
\label{sec:theorems}

\begin{theorem}[Fibrewise conservation]
\label{thm:conservation}
Fix a structural package $\mathfrak s$ and let
$f:\mathcal X\to\mathcal Y$ be an isomorphism in the fibre $\Car_{\mathfrak s}$.  Then
$\Dsharp_{\mathfrak s}(f)$ is an isomorphism in $\mathbf{Prof}_{\mathfrak s}$.  In particular:
\begin{enumerate}[label=(\roman*)]
\item a homotopy equivalence in the topological fibre preserves $\Dtop$ but need not preserve a
metric, a spectrum, a chosen field package or a chosen spin lift;
\item a bundle isomorphism over the base preserves the declared characteristic-class package;
\item an isometry preserves scalar and TT spectra with multiplicities, and a spin-isometry
preserves the Dirac spectrum as well.
\end{enumerate}
Thus conservation is not a global slogan: it is always conservation inside the fixed structural
category whose morphisms have been declared admissible.
\end{theorem}

\begin{proof}
Homology, cohomology and homotopy data are natural in the topological fibre; characteristic
classes are natural under pullback in the bundle fibre; and the relevant elliptic operators are
unitarily conjugate in the metric or spin-metric fibre.  The theorem is the fibrewise assembly
of these standard invariance statements.  Dropping any part of the structure moves to a
different fibre and is therefore a forgetting operation, not a counterexample to conservation.
\end{proof}

\begin{theorem}[Attributed erasure]
\label{thm:erasure}
Let $e:\mathcal X\to\mathcal Y$ be a morphism in $\Car$.  If at least one component of
$\Dsharp(e)$ fails to be an isomorphism---a reduced homology group changes, the cup-product
structure collapses, a characteristic class is lost, a stable homotopy class dies, or part of a
spectrum is truncated---then $e$ is not invertible in the corresponding structural category.
Consequently every change of the profile is \emph{attributed}: it identifies a concrete
forgetting morphism in $\Forg$ (Definition~\ref{def:forg}), and the lost package components
quantify what was erased.
\end{theorem}

\begin{proof}
Functors preserve isomorphisms; if $e$ had an inverse in the structural category, every induced
map on the profile would be an isomorphism, contradicting the hypothesis.  The generation
statement is Definition~\ref{def:forg}.

\end{proof}

\begin{definition}[Mapping-cone loss invariant]
\label{def:cone-loss}
For a map $f:X\to Y$, define its cone defect by the mapping cone
\[
C_f:=Y\cup_f CX.
\]
The degree-$k$ homological loss and torsion loss are
\[
E_k^{H}(f):=\operatorname{rank}\,\widetilde H_k(C_f;\Z),\qquad
E_k^{\mathrm{Tor}}(f):=\operatorname{Tor}\,\widetilde H_k(C_f;\Z).
\]
The exact topological erasure package is
\[
\mathbb E(f):=\big([C_*(C_f)],\ \widetilde H_*(C_f),\
\tau_W(f)\ \text{when the simple-homotopy torsion is defined}\big).
\]
Here $\tau_W$ records simple-homotopy loss invisible to ordinary homology in acyclic cases.
\end{definition}

\begin{proposition}[Cone criterion for zero loss]
\label{prop:cone-zero}
For maps of finite CW type: if $f:X\to Y$ is a homotopy equivalence, then $C_f$ is
contractible, so contractibility of the cone is a necessary condition for zero topological
erasure.  Conversely, if $C_f$ is contractible then $f$ is an integral homology equivalence;
if in addition $X$ and $Y$ are simply connected, $f$ is a homotopy equivalence.  The
simple-connectivity hypothesis cannot be dropped: for an acyclic complex $A$ with
$\pi_1(A)\neq0$, the map $A\to\{*\}$ has contractible cone (the suspension $\Sigma A$ is
simply connected and acyclic, hence contractible) without being an equivalence.  If
$C_f$ is acyclic but not contractible, rank-level and torsion-level losses vanish while a
simple-homotopy defect may remain.
\end{proposition}

\begin{proof}
The first statement is standard.  For the converse, the cofiber long exact sequence gives that
$\widetilde H_*(C_f)=0$ iff $f_*$ is an isomorphism on integral homology; in the simply
connected case Whitehead's theorem upgrades a homology equivalence to a homotopy equivalence.
The counterexample follows from Hurewicz and Whitehead applied to $\Sigma A$.  The final
statement follows from the distinction between acyclicity and contractibility, with Whitehead
torsion recording the residual simple-homotopy obstruction where applicable.
\end{proof}

\begin{proposition}[Generator formulae for forgetting morphisms]
\label{prop:forg-generators}
For the four generator types of $\Forg$ the cone-defect package has the following concrete
forms.
\begin{enumerate}[label=(\roman*)]
\item If $q:X\to X/A$ collapses a cofibration subspace $A\subset X$, then
\[
C_q\simeq \Sigma A,
\qquad
\widetilde H_k(C_q)\cong \widetilde H_{k-1}(A).
\]
The collapse cost is read from the shifted homology of the collapsed subcarrier.
\item If $j:X\setminus A\hookrightarrow X$ is deletion or restriction, then
\[
\widetilde H_k(C_j)\cong H_k(X,X\setminus A),
\]
and for good subspaces of oriented manifolds this relative group may be evaluated by Alexander
or Poincare-Lefschetz duality.
\item If coefficients change by a ring map $R\to S$, the safe defect is computed by derived change
of coefficients,
\[
C_*(X;R)\longrightarrow C_*(X;R)\otimes_R^{\mathbf L}S,
\]
and its mapping cone.  This prevents torsion from being silently discarded by a non-flat
coefficient change.
\item If spectral coarse-graining changes a cutoff $\Lambda$ to $\Lambda'<\Lambda$, then for a
sector $P$
\[
E^{P}_{\mathrm{spec}}(\Lambda,\Lambda')
=\dim\operatorname{im}(\Pi^{P}_{\le\Lambda}-\Pi^{P}_{\le\Lambda'}),
\]
with a smoothed version obtained by replacing sharp projectors by a test function
$\chi(P/\Lambda^2)$.
\end{enumerate}
Consequently OP2 is not a request for one scalar loss number but for the morphism-resolved
profile defect $\mathbb E(f)$.
\end{proposition}

\begin{proposition}[Sector Landauer bound]
\label{prop:sector-landauer}
If a reset operation at temperature $T$ sends $N$ mutually distinguishable logical sectors to
one standard sector, then its average work satisfies
\[
\langle W\rangle\ge k_B T\ln N.
\]
For a non-uniform prior this becomes $\langle W\rangle\ge k_BT\,\Delta H$, with $\Delta H$ the
Shannon entropy decrease of the sector distribution.  Higher-degree costs therefore have the
safe decomposition
\[
Q_T(f)\ge k_BT\,H(\text{identified sectors})+\Gamma_{\mathrm{geom}}(f),
\]
or, for $N_k$ equiprobable sectors,
\[
Q_k\ge k_BT\ln N_k+\Gamma_k^{\mathrm{geom}},
\]
where only the entropy term is a direct Landauer consequence and $\Gamma_k^{\mathrm{geom}}$ is
a protocol- and platform-dependent structural cost.
\end{proposition}

\begin{proposition}[The minimal example]
\label{prop:s0}
The identification $q:\Szero\to\{*\}$ kills $\rH_{0}$ ($\Z\to0$) and admits no homotopy
inverse; deleting a point destroys the structure; and no map exists from the nonempty
$\Szero$ to $\varnothing$.  The minimal distinction is erased only by an exhibited operator of
type (i)--(iii), never by continuous deformation.
\end{proposition}

\begin{proposition}[Non-monotonicity: identifications are not pure erasure]
\label{prop:nonmonotone}
Forgetting morphisms need not decrease total content: the contractible interval $[0,1]$ has
$\Dd\equiv0$, yet $[0,1]/(0\!\sim\!1)\simeq\Sone$ has $\Dd_{1}=1$.  ``Erasure'' and
``creation'' are therefore per-operation bookkeeping on the full profile, not global
monotonicity statements; every change must be computed, not assumed.
\end{proposition}

\begin{remark}[The precise slogan]
Theorems~\ref{thm:conservation}--\ref{thm:erasure} do \emph{not} say ``nothing in the world
ever disappears.''  They say something sharper and defensible: \emph{within a fixed category of
structure-preserving morphisms, the disappearance of an invariant is impossible; outside it,
the disappearance is attributable to a specific non-invertible operator with a computable
loss}.
\end{remark}

\begin{anchor}[Landauer's principle as the thermodynamic face of attributed erasure]
\label{anchor:landauer}
The logically irreversible erasure of one elementary distinction (one bit) in an environment at
temperature $T$ dissipates at least $k_{B}T\ln2$~\cite{Landauer1961,Bennett1982}; the bound has
been approached experimentally in colloidal and nanomagnetic
systems~\cite{Berut2012,Jun2014,Hong2016}.  The topological statement (erasure requires a
non-invertible operator) and the thermodynamic statement (logical irreversibility has a minimum
dissipation) are two faces of one discipline: nothing erases a distinction for free.  We cite
this as contact with established physics; the constant $k_{B}T\ln2$ is not derived from
topology, and extensions to higher-degree erasures are hypotheses (Sec.~\ref{sec:qk}).
\end{anchor}

\section{The Operational Layer: Postulates}
\label{sec:physical}

\begin{postulate}[Carrier realization]
\label{post:carrier}
A distinction is physically realized only when represented on a carrier equipped with metric
and causal structure.  Without scale there is no distance; without time there is no causal
relation; without causal relation there is no physically unfolded information; without
distinguishability there is no observable phenomenon:
\begin{equation}
\text{distinction}\;\to\;\text{carrier}\;\to\;\text{order/phase}\;\to\;\text{causal relation}
\;\to\;\text{reconstructible information}.
\label{eq:chain}
\end{equation}
\end{postulate}

\begin{postulate}[Resolution grain]
\label{post:grain}
Within the effective model, $\lP$ (and $\tP=\lP/c$) is treated as the lower boundary of
operationally reconstructible states: effective physical reconstruction assigns distinguishable
causal-geometric states only at spatial scales $\ge\lP$ and temporal scales $\ge\tP$.  Below
the grain the model assigns no operationally verifiable state---a boundary of the
reconstruction language within the effective theory, not a derived microscopic wall.
\end{postulate}

\begin{anchor}[Geometric support for Postulate~\ref{post:grain}]
In the companion Ricci-type flow analysis~\cite{BP2026ricci}, a reflecting barrier at $a=\lP$
defines a forward-invariant region with $|\Rm|^{2}\le C/\lP^{4}$, mirroring a self-adjoint
Wheeler--DeWitt endpoint on $a\in[\lP,\infty)$ with zero probability flux through $a=\lP$; the
same work proves that natural internal forcings cannot control the singular round mode, so the
grain enters as a \emph{boundary of the effective state space}, not as a smooth regulator.
Postulate~\ref{post:grain} adopts exactly this reading and inherits its falsifier: any
operationally confirmed sub-grain distinguishability refutes the postulate as stated
(Sec.~\ref{sec:falsifiability}).
\end{anchor}

\section{Graded Spectral Capacities on \texorpdfstring{$\Sthree_{R}$}{S3}}
\label{sec:capacity}

\subsection*{Counting conventions}
All capacities in this section are dimension counts of band-limited eigenspaces over the field
appropriate to the sector.  The scalar count is the complex eigenspace dimension; for real
fields it counts the corresponding real mode coordinates after the usual reality condition.
The Dirac count includes both positive and negative Dirac branches and counts complex spinor
components.  The TT count uses the transverse-traceless symmetric tensor harmonics of the positive rough
(Bochner) Laplacian $-\nabla^{2}$ in the convention of~\cite{RubinOrdonez1984}, as stated in
Theorem~\ref{thm:tt}; on the Einstein background $\Sthree_{R}$ the Lichnerowicz operator
differs from $-\nabla^{2}$ by a fixed curvature-proportional shift, which reindexes the cutoff
$M_{TT}$ by an $O(1)$ constant but leaves the multiplicities and the leading Weyl coefficient
unchanged.  These conventions are
part of the definition of $N_P$ and prevent the capacity formulae from being attacked as
notation-dependent claims.

\begin{definition}[Operational alphabet at grain $\varepsilon$]
For a closed Riemannian (spin) manifold and an elliptic operator $P$ of Laplace or Dirac type,
the band-limited space at grain $\varepsilon$ is spanned by the eigenmodes with
$\sqrt{|\lambda|}\le1/\varepsilon$ (Laplace type) or $|\mu|\le1/\varepsilon$ (Dirac type); the
\emph{mode capacity} $N^{(P)}_{\varepsilon}$ is its vector-space dimension: the number of independently
specifiable field modes of type $P$ resolvable at scale $\varepsilon$.  It is not yet a bit
capacity.  A bit capacity is obtained only after specifying a finite operational alphabet per
mode, or an energy/occupation cutoff for the associated Fock or code space.
\end{definition}

Write $x=R/\varepsilon$ throughout.

\begin{theorem}[Scalar capacity]
\label{thm:scalar}
On $\Sthree_{R}$ the scalar Laplacian has eigenvalues $\lambda^{(0)}_{n}=n(n+2)/R^{2}$ with
multiplicities $(n+1)^{2}$, $n\ge0$~\cite{PB2026susp}.  With
$M_{0}=\big\lfloor\sqrt{1+x^{2}}\big\rfloor$,
\begin{equation}
N_{0}(\varepsilon)=\sum_{n=0}^{M_{0}-1}(n+1)^{2}
=\frac{M_{0}(M_{0}+1)(2M_{0}+1)}{6}
\;\sim\;\frac{x^{3}}{3}.
\label{eq:N0}
\end{equation}
\end{theorem}

\begin{proof}
$n(n+2)\le x^{2}\iff(n+1)^{2}\le1+x^{2}\iff n+1\le M_{0}$; sum the multiplicities.
\end{proof}

\begin{theorem}[Transverse-traceless tensor capacity]
\label{thm:tt}
TT symmetric tensor harmonics on $\Sthree_{R}$ have rough-Laplacian eigenvalues
$\lambda^{(TT)}_{k}=\big(k(k+2)-2\big)/R^{2}$ with total multiplicity $2(k-1)(k+3)$ at each
$k\ge2$~\cite{RubinOrdonez1984,Sandberg1978}; in the labeling $n=k+1\ge3$
of~\cite{RubinOrdonez1984} these read $(n^{2}-3)/R^{2}$ with degeneracy $2(n^{2}-4)$.  With
$M_{TT}=\big\lfloor\sqrt{x^{2}+3}-1\big\rfloor$ (and $N_{TT}=0$ if $M_{TT}<2$),
\begin{equation}
N_{TT}(\varepsilon)=\sum_{k=2}^{M_{TT}}2(k-1)(k+3)
=\frac{M_{TT}(M_{TT}-1)(2M_{TT}+11)}{3}
\;\sim\;\frac{2x^{3}}{3}.
\label{eq:NTT}
\end{equation}
\end{theorem}

\begin{proof}
$k(k+2)-2\le x^{2}\iff(k+1)^{2}\le x^{2}+3\iff k\le\sqrt{x^{2}+3}-1$.  The closed form follows
by summing $2(k^{2}+2k-3)$ from $k=2$ to $M_{TT}$ and is verified by direct evaluation at
$M_{TT}=2,4,9,99$.
\end{proof}

\begin{theorem}[Dirac capacity]
\label{thm:dirac}
The Dirac operator on round $\Sthree_{R}$ has eigenvalues $\mu^{(D)}_{n}=\pm(n+\tfrac32)/R$
with multiplicity $(n+1)(n+2)$ for each sign, $n\ge0$~\cite{Baer1996,PB2026susp}.  With
$M_{D}=\big\lfloor x-\tfrac32\big\rfloor$ (and $N_{D}=0$ if $M_{D}<0$),
\begin{equation}
N_{D}(\varepsilon)=2\sum_{n=0}^{M_{D}}(n+1)(n+2)
=\frac{2(M_{D}+1)(M_{D}+2)(M_{D}+3)}{3}
\;\sim\;\frac{2x^{3}}{3}.
\label{eq:ND}
\end{equation}
\end{theorem}

\begin{proof}
$(n+\tfrac32)/R\le1/\varepsilon\iff n\le x-\tfrac32$; and
$\sum_{n=0}^{M}(n+1)(n+2)=(M+1)(M+2)(M+3)/3$.
\end{proof}

\begin{remark}[Weyl consistency and graded structure]
\label{rem:weyl}
The \emph{leading asymptotic terms} of \eqref{eq:N0}--\eqref{eq:ND} agree exactly with the
Weyl law for elliptic operators on closed $3$-manifolds and on vector bundles: for the scalar,
$N\sim\operatorname{Vol}\cdot\lambda^{3/2}/(6\pi^{2})$ with
$\operatorname{Vol}(\Sthree_{R})=2\pi^{2}R^{3}$ gives $x^{3}/3$; the TT and Dirac coefficients
are exactly twice the scalar one, matching the rank of the local field degrees of freedom (two
polarizations; the two signs of the Dirac spectrum counted as spinorial sectors).  The finite $N$'s themselves come from the explicit
$\Sthree$ spectra, not from the asymptotics.  Two consequences deserve emphasis: capacity is
naturally \emph{graded by field type} rather than a single number, and the thresholds differ
(Dirac and TT modes switch on later than scalar), which matters at moderate $x$.
\end{remark}

\begin{table}[ht]
\centering
\small
\begin{tabular}{r r r r}
\toprule
$x=R/\varepsilon$ & $N_{0}(x)$ scalar & $N_{TT}(x)$ TT tensor & $N_{D}(x)$ Dirac\\
\midrule
1   & 1        & 0        & 0\\
2   & 5        & 0        & 4\\
3   & 14       & 10       & 16\\
5   & 55       & 76       & 80\\
10  & 385      & 696      & 660\\
50  & 42\,925  & 85\,456  & 83\,300\\
100 & 338\,350 & 675\,906 & 666\,600\\
\bottomrule
\end{tabular}
\caption{Exact graded capacities on $\Sthree_{R}$ from
\eqref{eq:N0}--\eqref{eq:ND}.  The scalar count is one sector of the full graded capacity; the
TT and Dirac sectors activate at higher $x$ and dominate asymptotically with coefficient
$2/3$.}
\label{tab:capacity}
\end{table}

\begin{corollary}[Planck-grain capacity]
\label{cor:planck}
Under Postulate~\ref{post:grain} ($\varepsilon=\lP$), a quasi-local domain reconstructed as
$\Sthree_{R}$ carries the finite graded capacity
$\big(N_{0},N_{TT},N_{D}\big)$ of Theorems~\ref{thm:scalar}--\ref{thm:dirac} with $x=R/\lP$.
Finiteness is a theorem given the postulates; the identification $\varepsilon=\lP$ is the
postulate, not the theorem.  Whether this kinematic count is also the \emph{physical} count in
a gravitating domain is the open problem of Sec.~\ref{sec:tension}.
\end{corollary}

\section{Capacity Along Reconstruction Flow}
\label{sec:flow-capacity}

Sharp spectral counts jump when an eigenvalue crosses the cutoff.  A differentiable flow theory
therefore uses the heat-trace or smooth spectral-action capacity.

\begin{definition}[Heat-trace capacity]
For a smooth family of positive elliptic operators $P_\tau$ on a fixed carrier, define
\[
\mathcal N_P(s,\tau):=\operatorname{Tr}(e^{-sP_\tau}),
\qquad s=\varepsilon^2.
\]
Equivalently, with a smooth cutoff $\chi$, one may use
$S_\chi(\tau)=\operatorname{Tr}\chi(P_\tau/\Lambda^2)$.
\end{definition}

\begin{proposition}[Flow derivative]
\label{prop:heat-flow}
Under standard smoothness and trace-class hypotheses,
\[
\partial_\tau \mathcal N_P(s,\tau)
=-s\operatorname{Tr}\big((\partial_\tau P_\tau)e^{-sP_\tau}\big),
\]
where the formula follows from Duhamel's identity and cyclicity of the trace, so no commutation
of $P_\tau$ with $\partial_\tau P_\tau$ is assumed; and
\[
\frac{dS_\chi}{d\tau}
=\frac{1}{\Lambda^2}\operatorname{Tr}\big[\chi'(P_\tau/\Lambda^2)\partial_\tau P_\tau\big].
\]
For a round $\Sthree_{a(\tau)}$ with $x=a/\varepsilon$, the leading sharp-count asymptotics give
\[
\dot N_0\sim \frac{a^2\dot a}{\varepsilon^3},\qquad
\dot N_{TT}\sim \frac{2a^2\dot a}{\varepsilon^3},\qquad
\dot N_D\sim \frac{2a^2\dot a}{\varepsilon^3}.
\]
For the unnormalised round Ricci flow $\dot a=-2/a$, these become
\[
\dot N_0\sim -\frac{2a}{\varepsilon^3},\qquad
\dot N_{TT}\sim -\frac{4a}{\varepsilon^3},\qquad
\dot N_D\sim -\frac{4a}{\varepsilon^3}.
\]
Thus reconstruction flow erases operational distinguishability sector by sector; a Planck
barrier $a\ge\ell_P$ becomes a lower spectral latch rather than a verbal cutoff.
\end{proposition}

\section{Causal Integration: The Physical Integration Ratio}
\label{sec:pir}

Capacity is a property of the domain; \emph{integrated} information is bounded by causal
accessibility.  On the compact causal spacetime
$(\Sthree_{R}\times\R,\ g_{R}-c^{2}dt^{2})$ the causal diameter is $D_{c}=\pi R$ and full
causal coherence is reached at
\begin{equation}
\tcoh=\frac{\pi R}{c}
\end{equation}
\cite{PB2026light}.

\begin{lemma}[Geodesic ball volume on $\Sthree_{R}$]
\label{lem:ball}
The geodesic ball of radius $s\in[0,\pi R]$ has volume
\begin{equation}
V_{B}(s)=\int_{0}^{s}4\pi R^{2}\sin^{2}\!\frac{\rho}{R}\,d\rho
=2\pi R^{2}s-\pi R^{3}\sin\frac{2s}{R},
\qquad V_{B}(\pi R)=2\pi^{2}R^{3}=\operatorname{Vol}(\Sthree_{R}).
\end{equation}
\end{lemma}

\begin{definition}[Accessible fraction and PIR]
\label{def:pir}
The causally accessible volume fraction at time $T$ is
\begin{equation}
F(T)=
\begin{cases}
\dfrac{1}{\pi}\left(\dfrac{cT}{R}-\dfrac12\sin\dfrac{2cT}{R}\right), & 0\le T\le\tcoh,\\[1.4ex]
1, & T\ge\tcoh .
\end{cases}
\label{eq:F}
\end{equation}
For an operational alphabet of capacity $N_{\varepsilon}$, the \emph{Physical Integration
Ratio} is
\begin{equation}
\PIR(T):=\frac{N^{\mathrm{int}}_{\varepsilon}(T)}{N_{\varepsilon}},
\end{equation}
the fraction of the capacity causally integrable by an observer by time $T$.
\end{definition}

\begin{proposition}[PIR in the uniform-support model]
\label{prop:pir}
In the homogeneous saturation model---modes uniformly supported over the domain, integration
limited only by the light cone---$\PIR(T)=F(T)$.  Then:
\begin{enumerate}[label=(\alph*)]
\item \emph{cubic onset:}\quad
$\PIR(T)\sim\dfrac{2\pi^{2}}{3}\Big(\dfrac{T}{\tcoh}\Big)^{3}$ as $T\to0$
(locally Euclidean growth);
\item \emph{exact half-time symmetry:}\quad
$\PIR(\tcoh-T)=1-\PIR(T)$, in particular $\PIR(\tcoh/2)=\tfrac12$;
\item \emph{no early shortcut:}\quad no internal observer or computation integrates the full
Planck-grain capacity of Corollary~\ref{cor:planck} before $\tcoh$.
\end{enumerate}
\end{proposition}

\begin{proof}
(a) Expand \eqref{eq:F}: with $u=cT/R$, $\sin2u=2u-\tfrac43u^{3}+O(u^{5})$ gives
$F=\tfrac{2}{3\pi}u^{3}+O(u^{5})$; substitute $u=\pi T/\tcoh$.
(b) With $u'=\pi-u$: $F(u')=\tfrac1\pi(\pi-u+\tfrac12\sin2u)=1-F(u)$.
(c) $F(T)<1$ for $T<\tcoh$.

\end{proof}

\begin{proposition}[Shell-complete projector form of PIR]
\label{prop:shell-pir}
Let $P$ be a homogeneous elliptic operator on the round $\Sthree_R$, and let
$\Pi_{\Lambda}^{P}$ be the spectral projector onto a complete set of eigenshells below the
cutoff.  Then for every geodesic ball $B_s\subset \Sthree_R$,
\[
\frac{\operatorname{Tr}(1_{B_s}\Pi_{\Lambda}^{P}1_{B_s})}
{\operatorname{Tr}\Pi_{\Lambda}^{P}}
=
\frac{\operatorname{Vol}(B_s)}{\operatorname{Vol}(\Sthree_R)}.
\]
Thus, in the homogeneous shell-complete case, the volume formula for $\PIR$ is not merely a
uniform-support ansatz; it is the diagonal constancy of a symmetry-invariant spectral
projector.
\end{proposition}

\begin{proof}
The kernel of a shell-complete projector is invariant under the transitive isometry action of
$\Sthree_R$.  Therefore its diagonal value is constant.  Integrating this constant over
$B_s$ and over all of $\Sthree_R$ gives the stated ratio.
\end{proof}

\begin{remark}[Inhomogeneous extension]
For a general compact carrier or non-homogeneous domain the exact definition becomes
\[
\PIR_{P,\Lambda}(B)=
\frac{\int_B K_{\Lambda}^{P}(x,x)\,dV}{\int_M K_{\Lambda}^{P}(x,x)\,dV},
\]
where $K_{\Lambda}^{P}$ is the spectral-projector kernel.  Local Weyl asymptotics then give a
controlled high-cutoff approximation.  This closes the homogeneous part of OP5 and turns the
non-homogeneous part into a kernel-estimation problem.
\end{remark}

\begin{remark}[Mode-selected scalar correction]
For the unit round $\Sthree$ scalar zonal harmonic of degree $n$,
\[
\phi_n(\chi)=\frac{1}{\sqrt 2\pi}\frac{\sin((n+1)\chi)}{\sin\chi},
\]
the mass inside the geodesic ball of angular radius $\chi$ is
\[
w_n(\chi)=\int_{B_\chi}|\phi_n|^2dV
=\frac{\chi}{\pi}-\frac{\sin(2(n+1)\chi)}{2\pi(n+1)}.
\]
For $n=0$ this reduces to the volume fraction $F(\chi)$; for incomplete or mode-selected
alphabets the oscillatory term is the first explicit correction to shell-complete PIR.
\end{remark}

\begin{table}[ht]
\centering
\small
\begin{tabular}{r r}
\toprule
$\tau=T/\tcoh$ & $\PIR(\tau)$\\
\midrule
0.00 & 0.000000\\
0.10 & 0.006451\\
0.25 & 0.090845\\
0.50 & 0.500000\\
0.75 & 0.909155\\
0.90 & 0.993549\\
1.00 & 1.000000\\
\bottomrule
\end{tabular}
\caption{Causal integration on $\Sthree_{R}$ in the uniform-support model.  The onset is
cubic; the curve is exactly antisymmetric about the half-coherence time.}
\label{tab:pir}
\end{table}

\begin{definition}[Reconstructible information of a quasi-local domain]
\label{def:reconstructible}
The \emph{reconstructible information} of a quasi-local domain at time $T$ is the
distinguishability realized in the graded operational alphabet of
Sec.~\ref{sec:capacity}, restricted to the causally accessible region at $T$
(Definition~\ref{def:pir}), above the grain of Postulate~\ref{post:grain}.  This defines the
\emph{object}; given a source distribution on the alphabet, its \emph{quantity} is measured by
the standard Shannon functional.  The framework supplies the alphabet; Shannon's theory then
applies unchanged.
\end{definition}

\begin{remark}[Two limits of reconstructible information]
Extraction is locked between two boundaries.  \emph{Microscopic:} the grain excludes infinite
information densities; the sub-grain interval carries no weight in the reconstruction measure.
\emph{Macroscopic:} horizon formation, dilution, and expansion degrade the operational alphabet
and lengthen effective coherence, limiting the causally synchronizable fraction.  Within the
companion reconstruction-flow reading~\cite{BP2026ricci}, the past limit is the grain boundary
of smooth geometry and the future limit is loss of recoverable quasi-local
structure---\emph{insufficient distinguishability}, not contraction.
\end{remark}

\section{Higher Erasure Costs: The \texorpdfstring{$Q_{k}$}{Qk} Hypotheses}
\label{sec:qk}

Landauer's bound prices the erasure of a degree-$0$ distinction.  The graded profile of
Sec.~\ref{sec:profile} invites---but does not prove---a hierarchy of prices for erasing
higher-degree distinguishability.  Status discipline is essential here: the following are
\emph{hypotheses defining a research programme}, not theorems.

\begin{hypothesis}[$Q_{k}$ family]
\label{hyp:qk}
Let a physical system encode $\mathcal N_{k}$ distinguishable degree-$k$ sectors (components,
cycles, membranes/domains, or topological/fusion sectors) at temperature $T$.  Then the minimal
free-energy cost of a logically irreversible reset of the degree-$k$ distinction obeys
\begin{equation}
\Delta F_{k}\;\ge\;k_{B}T\ln\mathcal N_{k}\;+\;\Gamma_{k},
\end{equation}
where $\Gamma_{k}\ge0$ is the structural cost of destroying the corresponding carrier (loop
tension, surface energy, topological gap, defect energy, braiding/fusion reset work).  For
$k=0$, $\Gamma_{0}=0$ and $\mathcal N_{0}=2$ recover the Landauer bound.
\end{hypothesis}

\begin{center}
\small
\begin{tabularx}{\textwidth}{@{}lYYYl@{}}
\toprule
\textbf{Hyp.} & \textbf{What is erased} & \textbf{Minimal structural form} &
\textbf{Test platform} & \textbf{Status}\\
\midrule
$Q_{0}$ & binary distinction & $\Delta F_{0}\ge k_{B}T\ln2$ &
colloidal traps, nanomagnets~\cite{Berut2012,Jun2014,Hong2016} & established\\
$Q_{1}$ & cycle / phase loop & $\Delta F_{1}\gtrsim k_{B}T\ln\mathcal N_{1}+\tau_{\ell}L$ &
superconducting vortices, persistent-current rings & hypothesis\\
$Q_{2}$ & membrane / domain & $\Delta F_{2}\gtrsim k_{B}T\ln\mathcal N_{2}+\sigma A$ &
domain walls, interface reset & hypothesis\\
$Q_{\mathrm{top}}$ & anyon / fusion sector &
$\Delta F_{\mathrm{top}}\gtrsim k_{B}T\ln\mathcal D_{\mathrm{fus}}+n\Delta_{\mathrm{gap}}$ &
toric-code / surface-code, FQH, Majorana platforms~\cite{Kitaev2003} & hypothesis\\
\bottomrule
\end{tabularx}
\end{center}

\begin{remark}[How to test]
Surface-code and toric-code architectures encode logical information in homology classes of a
surface~\cite{Kitaev2003}; a forced reset of a nontrivial class to the trivial one is almost
literally the attributed-erasure operator of Theorem~\ref{thm:erasure} made measurable.  The
decisive experiments are calorimetric, not only logical: compare minimal dissipation of a
fusion-sector reset with single-bit erasure at matched temperature and cycle time; for vortex
and domain carriers, compare annihilation work with defect length/area.  In this form the
programme becomes a \emph{topological thermodynamics of erasure} rather than a philosophy of
information.
\end{remark}

\section{Quantum Information: Carrier-Resolved Resources}
\label{sec:quantum}

Normalized qubit states form $\Sthree\subset\C^{2}$; ray states form $\CP^{1}\simeq\Stwo$; the
projection is the Hopf map, whose nontriviality ($c_{1}=1$) is the exact topological reason why
global phase and Bloch direction do not factorize~\cite{PB2026susp}.  DRD does not replace
quantum capacities; it \emph{decomposes the carrier of quantum distinguishability}.  For a
$d$-level pure system the same carrier tower is
\[
\S^{2d-1}\longrightarrow \CP^{d-1},
\qquad H^*(\CP^{d-1};\Z)=\Z[\alpha]/(\alpha^d),
\qquad c_1(\mathcal O(-1))=-\alpha.
\]
For mixed states, the carrier is stratified by the spectrum of the density matrix; generic
coadjoint orbits are flag manifolds, while degeneracies give lower-dimensional strata.  Thus
\[
\Dsharp_Q(d)=\big(H^*(\CP^{d-1}),c_1(\mathcal O(-1)),[\omega_{FS}],
\Omega_{\mathrm{Berry}},\text{reference-frame data},\text{POVM alphabet}\big)
\]
is the safe quantum profile: topological obstruction plus operational measurement data, not a
replacement for quantum mechanics.

\begin{definition}[Carrier-resolved qubit resources]
\label{def:qubit}
For the qubit carrier we distinguish the triple
\begin{equation}
C^{\mathrm{DRD}}_{\mathrm{qubit}}
=\big(C_{\mathrm{proj}},\,C_{\mathrm{phase\text{-}ref}},\,C_{\mathrm{hol}}\big):
\end{equation}
$C_{\mathrm{proj}}$, distinguishability realized on the ray space $\Stwo$ (projective
measurements, no phase reference); $C_{\mathrm{phase\text{-}ref}}$, fibre-sensitive
distinguishability operational only relative to an external phase reference (interferometry);
$C_{\mathrm{hol}}$, the global holonomy/obstruction datum of the carrier ($c_{1}=1$), a
carrier-level obstruction rather than a per-state resource: it prices the impossibility of a
global phase convention.  In particular $C_{\mathrm{hol}}$ is not hidden per-state information
in an absolute phase; only relative phase with respect to an operational reference frame can
enter an accessible alphabet.
\end{definition}

\begin{remark}[Carrier-resolved reading of the Holevo bound (Interpretation)]
The Holevo bound $I_{\mathrm{acc}}\le\chi$~\cite{Holevo1973} is left unchanged.  What DRD adds
is an interpretation of \emph{which part} of $\chi$ is accessible in a given operational
setting: if the $\Sone$ fibre is quotiented by absolute phase equivalence, only
$C_{\mathrm{proj}}$-type distinguishability is operational; a shared phase reference promotes
part of the fibre information to the operational alphabet.  This is a conceptual refinement of
the alphabet, not a new theorem about channel capacity, and is labelled accordingly.
\end{remark}

\begin{remark}[Reference-frame discipline]
The phase-sensitive part of the quantum profile is operational only relative to a reference
frame.  For a symmetry group $G$, the no-reference-frame description is represented by
twirling,
\[
\Delta_G(\rho)=\int_G U_g\rho U_g^\dagger\,dg,
\]
and the accessible ensemble distinguishability is reduced from the usual Holevo quantity to the
corresponding twirled quantity.  Current quantum-reference-frame work makes this relational
reading explicit for general symmetries and for gauge systems, including general relativity
where coordinate frames themselves are gauge-dependent~\cite{Vanrietvelde2023,Thiemann2026}.
Crossed-product/QRF constructions further show how observer clocks and entropy definitions are
linked in semiclassical gravity~\cite{DeVuyst2025}.  DRD uses these results only to specify the
operational alphabet; it does not modify the Holevo theorem.
\end{remark}


\begin{remark}[Topological entanglement as carrier-level distinguishability (related work)]
Topological entanglement entropy~\cite{KitaevPreskill2006,LevinWen2006} isolates a universal,
deformation-insensitive constant in many-body entanglement---distinguishability stored at the
carrier level rather than in local data, in the spirit of the bundle layer of $\Dop$.
Multi-boundary entanglement in Chern--Simons theory on link complements in
$\Sthree$~\cite{Balasubramanian2017} ties such carrier-level entanglement to link topology.
We record these as adjacent results supporting the carrier viewpoint; no claim of derivation is
made in either direction.
\end{remark}

\section{Volume--Area Selection and Holonomy-Compatible Boundary Admissibility}
\label{sec:tension}

This point must not be talked away; it is made central.  The graded capacities of
Sec.~\ref{sec:capacity} count \emph{kinematically resolvable} band-limited modes and scale
with volume, $N\sim(R/\varepsilon)^{3}$.  Covariant entropy bounds scale with area,
$S\lesssim A/4\lP^{2}$~\cite{Bousso1999}; holographic entropy formulae and their quantum
extensions replace naive bulk counting by extremal or quantum-extremal generalized entropy
data~\cite{RyuTakayanagi2006,EngelhardtWall2015}.  The companion suspension paper adopts area
counting for its entropy hierarchy for exactly this reason~\cite{PB2026susp}.  DRD does not
take any boundary-projection ontology as primitive.  The internal geometry of the programme remains
Hopf/holonomy-based: the nontrivial phase--direction carrier is the Hopf fibration and its
obstruction/holonomy data, not a claim that the bulk is ontologically projected from a
boundary.  The fork is therefore principled but modest: either $N_{\varepsilon}$ counts all
kinematic modes while gravity selects a physically realizable subspace, or the operational
alphabet itself must be redefined in the presence of gravitational constraints.

\begin{tcolorbox}[colback=gray!4,colframe=gray!50,title=Holonomy discipline]
The word ``boundary'' in this section means an admissibility constraint or code-subspace
selection rule.  It is not an ontological replacement of the carrier by a boundary screen.
The native DRD structure remains holonomic/Hopf-theoretic: phase, spin, $c_1$, Berry/Hopf
holonomy and the carrier profile $\Dsharp$ are primary; area bounds only constrain which
parts of the kinematic alphabet can be physically realized.
\end{tcolorbox}

\begin{center}
\small
\begin{tabularx}{\textwidth}{@{}lYY@{}}
\toprule
\textbf{Model} & \textbf{Idea} & \textbf{What counts as physical}\\
\midrule
Gravitational dressing & local modes are not gauge-invariant without gravitational anchoring &
only relationally dressed observables\\
Boundary-code compression & kinematic alphabet restricted by boundary/area admissibility &
a selected code subspace of dimension $\lesssim A/\lP^{2}$\\
Mode decoupling / backreaction & part of the modes does not survive self-gravity &
dynamically stable modes below a collapse threshold\\
\bottomrule
\end{tabularx}

\end{center}

\begin{definition}[Mode counts, bit capacities, and gravitational selector]
For a field sector \(P\), define the raw mode count and its causally accessible part by
\begin{align}
\nkin^{(P)}(\varepsilon)&:=N_P(\varepsilon),\\
\ncausal^{(P)}(T)&:=\PIR_P(T)\,\nkin^{(P)}(\varepsilon).
\end{align}
These are numbers of accessible modes, not bits.  To convert modes into a bit capacity one must
specify an operational finite alphabet per mode.  Under a working normalization with
\(\Bmode^{(P)}\) bits per accessible mode,
\begin{equation}
\Bcausal^{(P)}(T;\Bmode):=\Bmode^{(P)}\,\ncausal^{(P)}(T).
\end{equation}
For a binary truncation \(\Bmode=1\); for a finite \(q\)-level alphabet
\(\Bmode=\log_2 q\).  For an untruncated bosonic Fock space no finite bit capacity is defined
by a finite mode count alone; an energy or occupation cutoff is part of the data.

The physical code-bit capacity is
\begin{equation}
\Bphys^{(P)}(R,\varepsilon,T)
=\sup_{\mathcal H_{\mathrm{code}}\subseteq \mathcal H_{\mathrm{op}}^{(P)}(\le\Lambda,T)}
\log_2\dim\mathcal H_{\mathrm{code}},
\end{equation}
where \(\mathcal H_{\mathrm{op}}^{(P)}(\le\Lambda,T)\) is the finite operational Hilbert/code
space obtained after spectral cutoff, causal restriction and the chosen finite-alphabet or
energy cutoff.  The maximisation is constrained by gauge invariance, gravitational dressing,
generalized entropy or area bounds, and a backreaction/collapse threshold.

The gravitational selector is therefore the bit-ratio
\begin{equation}
\Sigma_{\mathrm{grav}}^{(P)}:=
\frac{\Bphys^{(P)}}{\Bcausal^{(P)}}\quad
(0\le\Sigma_{\mathrm{grav}}^{(P)}\le1),
\end{equation}
not a ratio between a bit count and a raw vector-space dimension.  A safe upper-bound form is
\begin{equation}
\Bphys^{(P)}\le
\min\left\{
\Bcausal^{(P)},\ \frac{2\pi E R}{\hbar c\ln2},\ \frac{A}{4\ell_P^2\ln2}
\right\},
\label{eq:bit-safe-bound}
\end{equation}
where the second term is the energy--radius bound~\cite{Bekenstein1981} and the third is the
area/covariant-entropy admissibility cap~\cite{Bousso1999,RyuTakayanagi2006,EngelhardtWall2015}.  At the collapse threshold
\(E\sim Rc^4/(2G)\) these two external bounds have the same leading order,
\begin{equation}
\frac{2\pi E R}{\hbar c}\sim\frac{\pi R^2}{\ell_P^2}\sim \frac{A}{4\ell_P^2},
\end{equation}
so a Planck-grain volume mode count is not itself the physical bit count.

The boundary-compression ratio is correspondingly a bit-normalized quantity,
\begin{equation}
\BCR^{(P)}:=\frac{A/(4\ell_P^2\ln2)}{\Bcausal^{(P)}}.
\end{equation}
In the one-bit-per-mode working normalization and at \(\varepsilon=\ell_P\), \(\BCR^{(P)}\)
is of order \(\ell_P/R\) for large \(R\).  Thus the volume--area tension becomes a variational
selector problem rather than a verbal contradiction.
\end{definition}

\begin{proposition}[Bit-safe selector bound]
\label{prop:selector-bound}
Let $\mathcal A_{\mathrm{grav}}^{(P)}$ be the class of code subspaces satisfying: finite
operational alphabet or energy cutoff; the gauge/diffeomorphism constraints; gravitational or
relational dressing; the applicable generalized-entropy or area admissibility condition; and a
backreaction/collapse threshold.  Then every admissible code obeys
\[
\log_2\dim\mathcal H_{\mathrm{code}}\le
\min\left\{\Bcausal^{(P)},\frac{2\pi ER}{\hbar c\ln2},
\frac{A}{4\ell_P^2\ln2}\right\},
\]
and therefore the selector is a contraction on bit capacity:
\[
0\le \Sigma_{\mathrm{grav}}^{(P)}\le 1.
\]
\end{proposition}

\begin{proof}
The first inequality is containment of the code in the causally accessible operational
Hilbert/code space.  The second and third are the imposed energy--radius and area/generalized
entropy admissibility constraints.  Taking the supremum over
$\mathcal A_{\mathrm{grav}}^{(P)}$ gives the stated bound and the ratio definition of
$\Sigma_{\mathrm{grav}}^{(P)}$ gives the contraction property.
\end{proof}

\begin{remark}[No double counting of gauge charges]
In gauge theories with matter, a rigid matter symmetry, a constant asymptotic gauge
transformation, a bulk current expression and a boundary flux expression can describe the same
physical charge.  DRD therefore counts such a charge once.  It must be labelled as a bulk
Noether charge, a constraint expression, a Noether identity, or a surface charge, and two
expressions are identified when the field equations and boundary conditions make them equal on
shell.
\end{remark}

\begin{remark}[Algorithmic front end for the boundary-code branch]
If the boundary-code branch of $\Sigma_{\mathrm{grav}}$ is used, the relevant optimization is
not a slogan about screens but an entropy-cone constrained code-subspace problem.  Modern
holographic-entropy methods translate linear inequalities into contraction maps and partial
cubes, and recent reinforcement-learning searches use graph min-cuts to probe extreme rays of
the holographic entropy cone~\cite{Bao2025,Bao2025Complete,HeLeeOoguri2026}.  In DRD this
machinery supplies candidate admissibility constraints for $\mathcal H_{\mathrm{code}}$; it
does not replace the carrier profile $\Dsharp$.
\end{remark}

\begin{remark}[Observable-algebra reading]
The dressing requirement is not optional bookkeeping.  Compactly supported operators coupled
to gravity require gravitational dressing, and local commuting subalgebras cease to be the
right primitive notion of subsystem~\cite{DonnellyGiddings2016}.  Recent algebraic approaches
therefore study relationally dressed observables, gravitational splittings and type-II
operator-algebra structures as candidates for the quantum-gravity notion of localized
information~\cite{Giddings2025,AhmadJefferson2024,Faulkner2024}.  This is precisely the
mathematical habitat of $\Sigma_{\mathrm{grav}}$.
\end{remark}


\begin{tcolorbox}[colback=blue!3,colframe=blue!40!black,title=Central remaining physical node]
The strict mathematical and operational programme can be closed up to the selector
$\Sigma_{\mathrm{grav}}=\Bphys/\Bcausal$.  The open physical question is not whether the
kinematic spectral alphabet can be counted; it can.  The question is which finite-alphabet or
energy-cutoff code bits inside that alphabet remain gauge-invariant, holonomy-compatible,
backreaction-stable, and compatible with covariant entropy bounds.
\end{tcolorbox}

\begin{remark}[Status]
The three models are competitors, not conclusions.  This is the place where the framework can
fail or deepen: if all realistic realizations of capacity are gravitationally truncated to area
scaling and the volume count is never physically attained, the interpretation of
$N_{\varepsilon}$ as physical (rather than kinematic) capacity must be revised
(Sec.~\ref{sec:falsifiability}).
\end{remark}

\section{Related Definition-First Programmes and Conjectural Bridges}
\label{sec:related}

\paragraph{Information cohomology (adjacent, complementary).}
Baudot and Bennequin proved that Shannon entropy is, up to a multiplicative constant, the
unique nontrivial $1$-cocycle of an information cohomology built on categories of observables,
with the chain rule as the cocycle condition; deformed coefficients yield Tsallis
entropy~\cite{BaudotBennequin2015}.  This is a definition-first algebraicization of the
information \emph{measure}.  DRD is a definition-first algebraicization of the information
\emph{alphabet and carrier}.  The two programmes are complementary floors of the same tower:
information cohomology derives the functional that DRD's operational alphabets are measured by.
The bridge is the functor
\[
\Dsharp(\mathcal X)\longmapsto \mathbf{Obs}_{\varepsilon,T,\Sigma}(\mathcal X)
\longmapsto H^1(\mathbf{Obs}_{\varepsilon,T,\Sigma};\mathcal F),
\]
where objects of $\mathbf{Obs}_{\varepsilon,T,\Sigma}$ are finite spectral partitions, POVM
partitions or causal/dressed subalgebras after resolution, causal and selector restrictions, and
morphisms are coarse-grainings.  A carrier map $f:X\to Y$ induces a
pullback of observable categories whenever the structure permits it; entropy is then a
cocycle on the downstream observable category, not the primitive definition of information.

\paragraph{Reconstruction-flow thermodynamics (conjectural bridge).}
Perelman's entropy functional makes Ricci flow gradient-like~\cite{Perelman2002}; proposed
correspondences relate that functional to Fisher information and, further, to the Bohm quantum
potential~\cite{Carroll2006}.  If such a correspondence holds in a controlled setting, the
reconstruction flow of~\cite{BP2026ricci} would acquire a direct statistical-distinguishability
meaning, and OP3 of Sec.~\ref{sec:open} (capacity under reconstruction flow) would gain a
thermodynamic potential.  We record this strictly as a \emph{conjectural bridge}: no claim is
made here that quantum mechanics is derived from geometric flow.

\section{Interpretive Layer}
\label{sec:interpretation}

Everything in this section is labelled \emph{Interpretation}: a reading discipline, not
physics.

\begin{remark}[Non-emergence]
The theory does not assert the absolute emergence of information, space, or geometry.
$\Szero$ gives minimal distinguishability; suspension changes its carriers
(Theorem~\ref{thm:transport}); $\Sthree$ gives the minimal joint spin/Hopf carrier of the
stated package; the grain gives the lower metric-causal boundary of physical reconstruction
(Postulate~\ref{post:grain}).  Within the fixed category of invertible morphisms the invariant
package is neither created nor erased; across categories, every change is attributed to an
exhibited operator.  Cessation of a conditioned representation is not annihilation of the
invariant; it is release from a conditioned carrier.
\end{remark}

\begin{remark}[Two protections]
First, against self-being: no element of the chain---distinction, carrier, alphabet,
observer---is posited as self-existing; each is given by a construction (cf.\ the tetralemma
discipline~\cite{Garfield1995}).  Second, against substantializing dependence itself: having
refuted self-being, one may not erect ``Interdependence'' as a new substance.  In mathematical
register: empty of own-nature, not empty of dependent functioning.
\end{remark}

\begin{remark}[The axis]
Within the effective model, the Planck grain is the lower boundary of operationally
reconstructible states; in the programme's fixed formula: \emph{the Planck boundary is not the
beginning of the universe; it is the universal lower boundary of quasi-local geometric
reconstruction}~\cite{BP2026ricci}.  In informational terms: not a first bit, but the grain
below which the very notion of an operational bit is not assigned.
\end{remark}

\begin{tcolorbox}[colback=gray!4,colframe=gray!50,title=Core principle (Interpretation)]
Information is dependently reconstructible distinguishability: a functorial package of
invariants of relational structure (Definitions~\ref{def:dtop}--\ref{def:dop}), preserved
under structure-preserving re-description (Theorem~\ref{thm:conservation}), transported
without loss across carriers (Theorem~\ref{thm:transport}), destroyed only by exhibited
non-invertible operators at a computable and thermodynamically priced cost
(Theorem~\ref{thm:erasure}, Anchor~\ref{anchor:landauer}), physically realized only through
carrier, grain, and causal order (Postulates~\ref{post:carrier}--\ref{post:grain}), graded and
finite on compact domains (Theorems~\ref{thm:scalar}--\ref{thm:dirac}), and causally
integrable no faster than $\PIR(T)$ permits (Proposition~\ref{prop:pir}).  Not self-being, not
nothing: dependently reconstructible distinguishability passing through a chain of carriers.
\end{tcolorbox}

\section{Falsifiability}
\label{sec:falsifiability}

The mathematical layer is proved, not testable.  The operational and programmatic layers have
explicit failure modes.

\begin{enumerate}[label=(\arabic*)]
\item \emph{Mathematical failure modes.}  If any fibrewise functor
$\Dsharp_{\mathfrak s}$ fails on its declared morphism class; if the cone criterion or any
formula of Proposition~\ref{prop:forg-generators} fails on a standard counterexample; if the
spectral conventions of Sec.~\ref{sec:capacity} yield a wrong count; or if
Proposition~\ref{prop:shell-pir} fails for a shell-complete homogeneous sector---the strict
theory does not stand.  These are good risks: checkable and taste-independent.
\item \emph{Compact-topology searches.}  The framework does not predict that the universe is
$\Sthree$; it says that \emph{if} a quasi-local domain is effectively reconstructed as a
compact $\Sthree$ carrier, the graded spectral and causal consequences of
Secs.~\ref{sec:capacity}--\ref{sec:pir} follow.  Matched-circle and correlation searches
constrain compact topologies~\cite{Cornish1998,Planck2018}; null results bound the radius
without falsifying the mathematical core; future CMB experiments extend sensitivity.
\item \emph{Grain falsifier.}  Any stable, operationally confirmed distinguishability
\emph{below} the effective grain boundary refutes Postulate~\ref{post:grain} as stated, and
with it Corollary~\ref{cor:planck}.
\item \emph{Capacity-interpretation falsifier.}  If all realistic capacity realizations are
gravitationally truncated to area scaling (Sec.~\ref{sec:tension}), the interpretation of
$N_{\varepsilon}$ as physical capacity must be revised to kinematic capacity plus a selection
rule.
\item \emph{$Q_{k}$ tests.}  Calorimetric comparisons of topological-sector resets against
single-bit erasure (Sec.~\ref{sec:qk}) can confirm, bound, or refute each $Q_{k}$
independently.
\item \emph{Legacy ansatz, retired.}  The golden-ratio log-periodicity
$\omega=2\pi/\ln\golden$ of the wider programme was directly tested by us against DESI~DR2 BAO
and Pantheon$+$ with null results~\cite{BP2026null}; no statement in this paper depends on it,
and its current status is recorded as unsupported.  A companion null-status note records the
algebraic co-scaling condition $R_{N}=\lP\golden^{N}$, $t_{N}=\tP\golden^{N}$,
$R_{N}/t_{N}=c$.  Since $\tP=\lP/c$, this is a causal-admissibility identity, not a derivation
of the numerical speed of light.  DRD does not depend on this identity.
\end{enumerate}

\section{Conclusion, Publication Roadmap, and Open Problems}
\label{sec:open}

The paper set out to define information before using it.  The result is a three-layer
framework: a proved mathematical core (functorial profile, categorical conservation, attributed
erasure, suspension transport with Freudenthal stabilization), an operational layer with exact
graded capacities and an exact causal-integration law conditional on two explicit postulates,
and a programmatic/interpretive layer strictly quarantined from both.  The claimed contribution
is exactly this: a definition-first, theorem-based, falsifiability-aware account of what a
physically realizable alphabet of distinguishable states is, what preserves it, what destroys
it and at what price, and how much of it a compact quasi-local domain can carry and causally
integrate.

\paragraph{Publication roadmap.}
For refereed publication the framework separates naturally into a bundle of three papers, each
with one strong thesis: (1)~the mathematical core---carrier categories, $\Dsharp$,
conservation and attributed erasure; (2)~the spectral-causal core---graded capacities
$N_{0},N_{TT},N_{D}$, Weyl consistency, and $\PIR$; (3)~the physical
programme---carrier-resolved quantum resources, $Q_{k}$ hypotheses with experimental schemes,
and the volume--area/holonomy-compatible boundary-admissibility problem.  The present master document holds all three with a
uniform status discipline.

\paragraph{Status of the open nodes.}
OP1 is closed at the categorical level as the fibrewise family $\Dsharp_{\mathfrak s}$ with computable truncation
$\Dsharp_{q,\Lambda}$; remaining work concerns fibre-by-fibre examples, computations, and extensions; OP2 is closed at the topological level by the cone-defect package and
its generator formulae, and reduced at the physical level to sector Landauer plus
$\Gamma_{\mathrm{geom}}$; OP3 is made differentiable by heat-trace capacity; OP5 is closed in
the homogeneous shell-complete case by Proposition~\ref{prop:shell-pir} and reduced generally
to spectral-projector kernel estimates; OP6 and OP7 have precise carrier-to-observable
formulations.  OP4 is the unique remaining physical node: derive, compute or experimentally constrain
$\Sigma_{\mathrm{grav}}$ beyond the upper-bound discipline of Proposition~\ref{prop:selector-bound}.

\begin{enumerate}[label=\textbf{OP\arabic*.},leftmargin=2.2em]
\item \textbf{Fibrewise profile validation.}  Audit $\Dsharp_{\mathfrak s}$ fibre by fibre:
continuous maps, smooth maps, bundle maps, isometries and spin-isometries must never be mixed
without explicitly changing fibres.
\item \textbf{Computable cone defects.}  Implement $\mathbb E(f)$ for the four generators of
$\Forg$, including torsion and coefficient-change defects, and connect each nonzero component
to a candidate protocol cost $\Gamma_{\mathrm{geom}}$.
\item \textbf{Heat-trace flow capacity.}  Replace sharp $N_P(\Lambda,\tau)$ by
$\mathcal N_P(s,\tau)$ for reconstruction flows; for Berger $\Sthree$, compute the branch
crossings by $SU(2)$ representation theory and compare with the round asymptotics of
Proposition~\ref{prop:heat-flow}.
\item \textbf{Gravitational selector.}  Define the bit-safe selector
$\Sigma_{\mathrm{grav}}=\Bphys/\Bcausal$ through gauge-invariant dressed observables,
boundary-code compression subspaces, generalized entropy, the energy--radius bound and
backreaction thresholds.  This is the central physical
problem of the programme.
\item \textbf{Mode-resolved PIR.}  Extend the shell-complete theorem to incomplete scalar, TT
and Dirac alphabets by explicit projector kernels, beginning with the zonal scalar formula.
\item \textbf{Quantum carrier profile.}  Develop the Hopf/flag/reference-frame profile
$\Dsharp_Q$ for pure and mixed states and make its relation to the Holevo quantity quantitative
without claiming a new capacity bound.
\item \textbf{Information-cohomology interface.}  Build the functor
$\Dsharp(\mathcal X)\to\mathbf{Obs}_{\varepsilon,T}(\mathcal X)$ and evaluate entropy cocycles
only after the operational alphabet has been reconstructed.
\end{enumerate}

The framework may fail; but if it fails, it fails by calculation---at a functoriality check,
at a capacity formula, at a calorimetric bound, at an observational anchor---not by taste.

\begin{thebibliography}{99}

\bibitem{Shannon1948}
C.~E.~Shannon, \emph{A mathematical theory of communication},
Bell Syst.\ Tech.\ J.\ \textbf{27}, 379--423, 623--656 (1948).

\bibitem{Kolmogorov1965}
A.~N.~Kolmogorov, \emph{Three approaches to the quantitative definition of information},
Problems Inform.\ Transmission \textbf{1}, 1--7 (1965).

\bibitem{Holevo1973}
A.~S.~Holevo, \emph{Bounds for the quantity of information transmitted by a quantum
communication channel}, Probl.\ Peredachi Inf.\ \textbf{9}, 3--11 (1973).

\bibitem{Landauer1961}
R.~Landauer, \emph{Irreversibility and heat generation in the computing process},
IBM J.\ Res.\ Dev.\ \textbf{5}, 183--191 (1961).

\bibitem{Landauer1991}
R.~Landauer, \emph{Information is physical}, Phys.\ Today \textbf{44}(5), 23--29 (1991).

\bibitem{Bennett1982}
C.~H.~Bennett, \emph{The thermodynamics of computation---a review},
Int.\ J.\ Theor.\ Phys.\ \textbf{21}, 905--940 (1982).

\bibitem{Berut2012}
A.~B\'erut, A.~Arakelyan, A.~Petrosyan, S.~Ciliberto, R.~Dillenschneider, and E.~Lutz,
\emph{Experimental verification of Landauer's principle linking information and thermodynamics},
Nature \textbf{483}, 187--189 (2012).

\bibitem{Jun2014}
Y.~Jun, M.~Gavrilov, and J.~Bechhoefer,
\emph{High-precision test of Landauer's principle in a feedback trap},
Phys.\ Rev.\ Lett.\ \textbf{113}, 190601 (2014).

\bibitem{Hong2016}
J.~Hong, B.~Lambson, S.~Dhuey, and J.~Bokor,
\emph{Experimental test of Landauer's principle in single-bit operations on nanomagnetic
memory bits}, Sci.\ Adv.\ \textbf{2}, e1501492 (2016).

\bibitem{Wheeler1990}
J.~A.~Wheeler, \emph{Information, physics, quantum: the search for links}, in
\emph{Complexity, Entropy, and the Physics of Information}, ed.\ W.~H.~Zurek
(Addison-Wesley, 1990).

\bibitem{Hatcher2002}
A.~Hatcher, \emph{Algebraic Topology} (Cambridge University Press, 2002).

\bibitem{MilnorStasheff}
J.~Milnor and J.~Stasheff, \emph{Characteristic Classes}
(Princeton University Press, 1974).

\bibitem{Baer1996}
C.~B\"ar, \emph{The Dirac operator on space forms of positive curvature},
J.\ Math.\ Soc.\ Japan \textbf{48}, 69--83 (1996).

\bibitem{RubinOrdonez1984}
M.~A.~Rubin and C.~R.~Ord\'o\~nez,
\emph{Eigenvalues and degeneracies for $n$-dimensional tensor spherical harmonics},
J.\ Math.\ Phys.\ \textbf{25}, 2888--2894 (1984).

\bibitem{Sandberg1978}
V.~D.~Sandberg, \emph{Tensor spherical harmonics on $S^{2}$ and $S^{3}$ as eigenvalue
problems}, J.\ Math.\ Phys.\ \textbf{19}, 2441--2446 (1978).

\bibitem{Kitaev2003}
A.~Yu.~Kitaev, \emph{Fault-tolerant quantum computation by anyons},
Ann.\ Phys.\ \textbf{303}, 2--30 (2003).

\bibitem{KitaevPreskill2006}
A.~Kitaev and J.~Preskill, \emph{Topological entanglement entropy},
Phys.\ Rev.\ Lett.\ \textbf{96}, 110404 (2006).

\bibitem{LevinWen2006}
M.~Levin and X.-G.~Wen, \emph{Detecting topological order in a ground state wave function},
Phys.\ Rev.\ Lett.\ \textbf{96}, 110405 (2006).


\bibitem{Balasubramanian2017}
V.~Balasubramanian, J.~R.~Fliss, R.~G.~Leigh, and O.~Parrikar,
\emph{Multi-boundary entanglement in Chern--Simons theory and link invariants},
JHEP \textbf{04}, 061 (2017).

\bibitem{Bekenstein1981}
J.~D.~Bekenstein, \emph{Universal upper bound on the entropy-to-energy ratio for bounded systems},
Phys.\ Rev.\ D \textbf{23}, 287--298 (1981).

\bibitem{Bousso1999}
R.~Bousso, \emph{A covariant entropy conjecture}, JHEP \textbf{07}, 004 (1999).

\bibitem{BaudotBennequin2015}
P.~Baudot and D.~Bennequin, \emph{The homological nature of entropy},
Entropy \textbf{17}, 3253--3318 (2015).

\bibitem{Perelman2002}
G.~Perelman, \emph{The entropy formula for the Ricci flow and its geometric applications},
arXiv:math/0211159 (2002).

\bibitem{Carroll2006}
R.~W.~Carroll, \emph{Fluctuations, Information, Gravity and the Quantum Potential}
(Springer, 2006).

\bibitem{Cornish1998}
N.~J.~Cornish, D.~N.~Spergel, and G.~D.~Starkman,
\emph{Circles in the sky: finding topology with the microwave background radiation},
Class.\ Quantum Grav.\ \textbf{15}, 2657--2670 (1998).

\bibitem{Planck2018}
Planck Collaboration, \emph{Planck 2018 results.\ VI.\ Cosmological parameters},
Astron.\ Astrophys.\ \textbf{641}, A6 (2020).


\bibitem{RyuTakayanagi2006}
S.~Ryu and T.~Takayanagi, \emph{Holographic derivation of entanglement entropy from AdS/CFT},
Phys.\ Rev.\ Lett.\ \textbf{96}, 181602 (2006).

\bibitem{EngelhardtWall2015}
N.~Engelhardt and A.~C.~Wall, \emph{Quantum extremal surfaces: holographic entanglement
entropy beyond the classical regime}, JHEP \textbf{01}, 073 (2015).

\bibitem{DonnellyGiddings2016}
W.~Donnelly and S.~B.~Giddings, \emph{Observables, gravitational dressing, and obstructions to
locality and subsystems}, Phys.\ Rev.\ D \textbf{94}, 104038 (2016).

\bibitem{Giddings2025}
S.~B.~Giddings, \emph{Quantum gravity observables: observation, algebras, and mathematical
structure}, arXiv:2505.22708 (2025).

\bibitem{AhmadJefferson2024}
S.~Ali~Ahmad and R.~Jefferson, \emph{Crossed product algebras and generalized entropy for
subregions}, SciPost Phys.\ Core \textbf{7}, 020 (2024).

\bibitem{Faulkner2024}
T.~Faulkner and A.~J.~Speranza, \emph{Gravitational algebras and the generalized second law},
JHEP \textbf{11}, 099 (2024).

\bibitem{Bao2025}
N.~Bao, N.~Cheng, S.~Hernandez-Cuenca, and V.~P.~Su, \emph{Towards a complete classification
of holographic entropy inequalities}, JHEP \textbf{03}, 117 (2025).

\bibitem{Bao2025Complete}
N.~Bao, N.~Cheng, S.~Hernandez-Cuenca, and V.~P.~Su, \emph{On the completeness of contraction
map proof method for holographic entropy inequalities}, JHEP \textbf{12}, 140 (2025).

\bibitem{HeLeeOoguri2026}
T.~He, J.~Lee, and H.~Ooguri, \emph{Exploring the holographic entropy cone via reinforcement
learning}, arXiv:2601.19979 (2026).

\bibitem{Vanrietvelde2023}
A.~Vanrietvelde, P.~A.~Hoehn, F.~Giacomini, and E.~Castro-Ruiz, \emph{Switching quantum
reference frames in the $N$-body problem and the absence of global relational perspectives},
Quantum \textbf{7}, 1088 (2023).

\bibitem{Thiemann2026}
T.~Thiemann, \emph{(Quantum) reference frames, relational observables, gauge reduction and
physical interpretation}, arXiv:2603.04072 (2026).

\bibitem{DeVuyst2025}
J.~De Vuyst, S.~Eccles, P.~A.~H\"ohn, and J.~Kirklin, \emph{Crossed products and
quantum reference frames: on the observer-dependence of gravitational entropy}, JHEP
\textbf{07}, 063 (2025).


\bibitem{PB2026susp}
A.~Preece and B.~Batenin, \emph{Suspension of Spheres and the Minimal Role of $\Sthree$ in
Spin, Hopf Geometry, and Irreversible Hierarchy}, research manuscript (2026).

\bibitem{BP2026ricci}
B.~Batenin and A.~Preece, \emph{Hopf-Lifted Ricci-Type Flow on $\Sthree$: A No-Go for Bounded
and Engineered Forcing, and a Planck-Barrier Invariant Region}, research manuscript (2026).

\bibitem{PB2026light}
A.~Preece and B.~Batenin, \emph{Light as Causal Coherence Structure in Compact $\Sthree$
Geometry}, research manuscript (2026).

\bibitem{BP2026null}
B.~Batenin and A.~Preece, \emph{A Null Test of Golden-Ratio Log-Periodicity
($\omega=2\pi/\ln\varphi$) against DESI DR2 BAO and Pantheon$+$}, research manuscript (2026).

\bibitem{Garfield1995}
N\=ag\=arjuna, \emph{M\=ulamadhyamakak\=arik\=a}; English translation and commentary:
J.~L.~Garfield, \emph{The Fundamental Wisdom of the Middle Way} (Oxford University Press,
1995).

\end{thebibliography}

\end{document}
